\begingroup
\setlength{\tabcolsep}{4pt}
\begin{longtable}{L{0.23\textwidth}L{0.30\textwidth}L{0.20\textwidth}L{0.17\textwidth}}
\toprule
Reprezentáció & Lényege & Előny & Hátrány \\
\midrule
\endfirsthead
\toprule
Reprezentáció & Lényege & Előny & Hátrány \\
\midrule
\endhead
Szomszédsági lista & Minden csúcshoz eltároljuk a belőle elérhető szomszédokat, súlyozott gráfnál a súlyokat is. & Ritka gráfnál tárhatékony, a BFS/DFS természetes bemenete. & Egy konkrét él létezésének ellenőrzése lassabb lehet. \\
Szomszédsági mátrix & $n\times n$ mátrix, ahol a cella az él létezését vagy súlyát adja meg. & Él létezése gyorsan ellenőrizhető, Floyd--Warshallhoz természetes. & $\Theta(|V|^2)$ tár akkor is, ha kevés él van. \\
Éllista & Az éleket külön listában tároljuk, például $(u,v,w)$ hármasokként. & Kruskal és Bellman--Ford esetén kényelmes. & Egy csúcs szomszédainak kigyűjtése külön index nélkül lassú lehet. \\
\bottomrule
\end{longtable}
\endgroup
